%!TeX root = ../incremental_SS_Translation.tex

\chapter{Śārīrasthāna 4:  On the Formation of the Embryo}


\section{Introduction}

This chapter opens with sixteen passages that discuss of the seven
of \sepl{tvac}{skin} and \sepl{kalā}{membrane} that form early
in the fetus's life.\footnote{On the system of the \dev{kalā} see
    \volcite{1}[183--184]{josi-maha}, \cite[227--228]{gupt-1983},
    \cite[6]{kutu-1962}, \volcite{1}[247--248 and notes]{meul-hist}.} This
    system of dermal and interstitial \sepl{kalā}{membrane} was not
    known to the \CS\ as such.  Rather, the \CS\ mentioned six kinds of
    \dev{tvac}\sse{tvac}{skin} with different names and characteristics.
    These were classified not according to appearence, as in the \SS, but
    mainly according to the diseases that they
    supported.\footnote{\label{sa4:kalā}\CS\ on \Ca{4.7.4}{337}. This
        contradiction between the \CS\ and the \SS\ was discussed by the
        commentator Cakrapāṇidatta (\emph{idem}).} The concept of interstitial
        skins was used in the \SS\ as an explanatory mechanism for the stages
        of snake envenomation.\footnote{See p.\,\pageref{ka4:kalā-vega}.} The
            \SS's concept of seven skins dominates the narrative of later
            works.\footnote{For example, the fourteenth-century
                \emph{Śārṅgadharasaṃhitā} (1.5.60, \cite[40]{sast-1931}) and the
                sixteenth-century \emph{Bhāvaprakāśa} (1.3.220--222,
                \volcite{1}[49]{brah-1935}), which gives seven skins, like the \SS,
                but names and describes them in the manner of the \CS. The \AS\
                (Śārīrasthāna 5.18, \cite[296--297]{atha-1980} gives both views.}

The chapter then moves to the description of the flow of substances in
the body and the conduits through which such flow takes place, and
then the formation of “organs.” \citet{zimm-1978,zimm-1978b} long ago
elegantly problematized the early twentieth-century attempts to equate
Ayurvedic descriptions and models in the light of post-Vesalian
European anatomical understandings. These efforts have only grown
stronger and less sophisticated with the continuing identification of
Ayurvedic medicine as a component of Indian national identity.
Although we use the language of “organs” in this translation, we do so
advisedly.  As \citet{zimm-1979} also argued, in Ayurveda, as in
pre-Aristotelian Greek medical thought, we have “Pathology without
Anatomy.”  That is to say, in Ayurveda, we have not so much 
anatomical descriptions of organs as individual physical objects, but
rather “a conceptual system of fluids, channels and diseases: fluids
and channels that symbolize organic functions, and diseases that
introduce imbalance and strain.”

\section{Literature} 

Meulenbeld offered an annotated overview of this chapter and a
bibliography of earlier scholarship to 2002 and, in his notes, citations
of the parallel passages in the \CS.\fvolcite{IA}[247--249]{meul-hist}   

The first draft of the following translation was prepared by Jan
Gerris, and later revised in collaboration Philipp Maas, John Nemec,
Dominik Wujastyk, Madhusudan Rimal, Kenneth Zysk and Andrey Klebanov.

\newpage

\section{Translation}

\begin{translation}

    % First draft by Jan Gerris, 2026-01-17
  
\bigskip       

\item[1]

Next we shall discuss the anatomical chapter on the analysis of the foetus.
 
\item[3]

\newenvironment{tightquote}
    {\list{}{\listparindent 1.5em%
          \itemindent    \listparindent
          \rightmargin   \leftmargin
          \parsep        \z@ \@plus\p@}%
      \item\relax}
  {\endlist}

Fire, and \se{soma}{liquid}, air, sattva, rajas, and tamas, the five
senses, the \se{bhūtātman}{elemental self}, and the mind
are the \sepl{prāṇa}{life principle}.\footnote{\dev{prāṇa} is here used to 
refer to all the components of a living being, not merely the five breaths.
On this passage, and its concept of multiple \dev{prāṇa}, see 
\cite[\S3.2.3]{kleb-2021b}.
On the early history of \dev{prāṇa}, see \cite{zysk-1993,zysk-2007}.
On the expression \dev{agnīṣoma},  “fire and liquid,” see 
\cite{wuja-2004,ange-2021}.

The word “\se{manas}{mind}” is present in the Nepalese version, but not
in the vulgate text.  The commentator Gayadāsa (fl.\ ca.\ 1000)
discussed this term, confirming that the word was present in the \SS\
text available to him, and he noted that his predecessor Jejjaṭa did
not think this was proper (\MS{Cambridge Add.2491}, f.\,33r):
\begin{quote}
\dev{jaḍastu bhūtātmaśabdena 
na mano'bhidhatte/ manograhaṇaṃ nādhīyate/  tanna
bhūtātmamanasorvyaktabhedāt/}\par  But Jejjaṭa does not intend “mind” by
the term “elemental self.”  The mention of the word “mind” is not
taught there. That is because there is an obvious difference between
the elemental self and the mind.
\end{quote}
This suggests that the word \emph{manas} was present in the earliest
\SS\ but was dropped as a result of Jejjaṭa's objection (that then
influenced Candraṭa's revision).}


\subsection{The seven skins}

\item[4]

You see, sperm and menstrual blood, maturing, get seven
\sepl{tvac}{skin}, for it is like the \sepl{santānika}{skin} on heated
milk.\footnote{If we take \dev{śukraśoṇitasyābhipacyamānasya} as a
    genitive absolute, we could read this statement as “even though [the
    admixture of] sperm and menstrual blood is forming, seven skins come
    into existence, like the skins on milk.”
    
 The following characterization of these skins is longer in the
vulgate text because it imports the concept of
“\sepl{adhiṣṭhāna}{foundation}.”  This concept of skins as
“foundations” is present in the \CS\ account of six skins, which are
described as being the foundations of various illnesses. See
pp.\,\pageref{ka4:kalā}, \pageref{sa4:kalā}.  It seems likely that
the vulgate here has been supplemented with material from the \CS.}

The first  of these, called  “\se{avabhāsinī}{Shining},”  makes all
colours shine and makes  visible five kinds of \se{chāyā}{complexion}.
It is the size of one-eighteenth part of a rice
grain.\footnote{\Dalhana{3.4.4}{350} interpreted \dev{vrīhi} “grain of
    rice” as being a barley corn (a standard unit of measurement).}

The second is called “\se{lohitā}{Red},” the size of one-sixteenth.
  
The third is “\se{śvetā}{White},” the size of one-twelfth. 
 
The fourth is “\se{tāmra}{Coppery},” one-eighth in size. 
 
The fifth is named “\se{vedanī}{Feeling},” one-fifth in size. 
 
The sixth is named “\se{rohiṇī}{Scarlet},” the size of one rice grain.

The seventh is named  “\se{śukradharā}{Semen-supporter},”  the size of
two rice grains.
 
Since it will be said, in the chapter about the belly, 
\begin{quote}
     Using the \se{vrīhimukha}{rice-tip instrument}, one penetrates
the measure of the thickness of a thumb or a a\se{aṅgula}{finger's
    breadth}.\footnote{This sentence is cited from \Su{4.14.18}{461}
    in the chapter calle \emph{Udarāṇāṃ cikitsitam} “On the therapy
    of the abdominal ailments.”  The measure
    \dev{aṅguṣṭhodarapramāṇa} “the size of the belly of a thumb” is
    \dev{viṃśatitamabhāgonaṣaḍyavapramāṇam} “one twentieth less than
    six barleycorns” (5$\frac{19}{20}$) according to
    \Dalhana{4.4.4}{355}. Cf.\ \cite[10]{josi-maha}.  The procedure
    described is the puncturing the abdomen and the insertion of a
    tube to release an abnormal buildup of fluid (ascites); it is
    today called paracentesis.  The citation from the
    \emph{Cikitsāsthāna} about piercing for ascites illustrates that
    the sum of the thicknesses of the above-mentioned layers
    corresponds to the thickness of the belly that has to be pierced.
    
    The “rice-tip” instrument is one of the surgical knives described
in \SS\ \Su{1.8}{35--41}, \cite[tr.][83--86]{wuja-2003}; see
\volcite{1}[257--261]{mukh-1913}. See
Fig.\,\ref{fig:vrihimukha-from-ss-1938} and the artists'
reconstructions at \volcite{2}[plate LXXI]{mukh-1913}.  }
 \end{quote}
    
     \begin{figure}
        \centering
        \includegraphics[width=0.7\linewidth]{"media/vrihimukha from SS 1938"}
        \caption{The \emph{Vrīhimukha} instrument, as illustrated in 
        \cite[37]{vulgate}.}
        \label{fig:vrihimukha-from-ss-1938}
    \end{figure}

 
 \subsection{The seven membranes}
 
 \item[5]
 
 You see, the seven \sepl{kalā}{membrane} also arise,  which are the 
 boundaries between the \sepl{dhātu}{tissue} and inner 
 \sepl{āśaya}{receptacle}.
 
 \item[6]
 There are two verses.
 \begin{sloka}
Just as one sees the \se{sāra}{core} when wood is being cut, 
 in the same way, one sees \se{dhātu}{tissue} when flesh is being 
 cut.\footnote{\dev{sāra} may here mean “pith,” i.e., the woody core of a 
 tree.}
\end{sloka}
 
\item[7]
 
 \begin{sloka}
And experts know that the \se{kalā}{membrane} parts are hidden by
\sepl{snāyu}{sinew} and covered by \se{jarāyu}{amnion} and also
surrounded by \se{śleṣman}{mucus}.\footnote{The series of seven \dev{kalā} 
membranes that are described in the following passages bear a similarity to the 
bodily constituents \dev{rasa}/\dev{tvak}--\dev{śukra} that are described in 
various places in the \SS\ and the \CS; see discussion by \citet{maas-2008}.}
\end{sloka}
 
 \item[8]
 
The first of these is called “\se{māṃsadharā}{flesh-support}.” 
 
From it arise ever new \sepl{pratāna}{ramification} in the flesh:
\sepl{sirā}{duct}, \sepl{snāyu}{sinew}, \sepl{dhamanī}{pipe}, and
\sepl{srotas}{tube}.\footnote{On these different conduits, see
    \cites[404--406]{wuja-2022}[xlvi--xlvii]{wuja-2003}[26--28]{ray-1980} and 
    the descriptions by 
    \tvolcite{2}[344--352]{dasg-medi}. The translation
    “pipe” for \dev{dhamanī} (from the \root\,\dev{dham} “blow”) is
    intended to suggest the primary function of transporting air;
    “vessel” would be an alternative translation. Adhyāyas 
    7 and 8 of the \emph{Śārīrasthāna} describe the \dev{sirā} and adhyāya 9 
    the \dev{dhamanī}, with the \dev{srotas} being described at the end of 
    chapter 9.}
  
\item[9]

On this there is the following:

\begin{sloka}
Conduits such as the ducts grow in the flesh just like lotus roots located in 
muddy water grow in all directions in the ground.
\end{sloka}

 \item[10]
 
The second one is called “\se{raktadharā}{blood-support}.” It is inside the 
flesh. 
From it arises blood, particularly in the \sepl{sirā}{duct} and in the
liver and spleen.
 
\item[11]

On this there is the following:

\begin{sloka}
    Just as milky sap of a sappy tree can flow from a tree that has been cut, 
    in the same way blood issues rapidly from injured flesh.
\end{sloka}

 % MS H, f. 167
 
 
 \item[12]
 
The third one is called named  \se{medodharā}{fat-support}.  You
see, all beings have fat in the abdomen and also marrow inside the large 
bones.\footnote{This view differs from that of the vulgate and Ḍalhaṇa who 
saw fat in smaller bones and marrow in larger ones.}
 
 \item[13]
 
 And there is a verse on this:
 
\begin{sloka}
Marrow is located especially inside in the thick bones.  Whereas, in
all others, fat is said to occur together with blood.\footnote{The
    vulgate text adds half a verse here introducing \dev{vasā} as a type
    of fat that is specifically present in “pure flesh.”  It seems that
    the Nepalese version does not use the term \sanengdev{vasā}{fat} at this 
    point.
    \Dalhana{3.4.13}{356} called \dev{vasā} a 
    \sanengdev{upadhātu}{subsidiary
    tissue}.   The term occurs more frequently in the \CS; in the \SS, it normally 
    refers to the fat of animals that live in wet, marshy environments.}
\end{sloka} 


\item[14]

The fourth one is called  \se{śleṣmadharā}{phlegm-support}, which is present 
in all the joints of all living beings.

 

\item[15]

There is a verse on this:

\begin{sloka}
Just as a wheel rotates smoothly on a lubricated axle, so the joints
move smoothly when combined with phlegm. 
\end{sloka} 
 
 
\item[16]

The fifth one, which is called the \se{purīṣadharā}{feces-support},
supports the four kinds  of food that have descended from the stomach
into the \se{pakvāśaya}{colon}.\footnote{Cf.\ \cite[177]{das-2003}.}
    Being located in the colon, it divides up the \se{koṣṭha}{abdominal
        cavity}.
 
 
 
 \item[17]
 
 There is a verse on this:
 \begin{sloka}     
      The impurity located at the \se{uṇḍuka}{lower intestine} is
divided by the \se{kalā}{membrane} called
\se{maladharā}{impurity-support}.\footnote{ \label{uṇḍuka}The
    term \dev{uṇḍuka}, “intestine, caecum, rectum or anal canal,” is
    of uncertain interpretation. Ḍalhaṇa said that the
    ordinary-language term for \dev{uṇḍuka} was \dev{poṭṭalakā}. The
    \SS\ defined \dev{uṇḍuka} as one of the main internal components
    of the body present in the trunk: “The trunk of the body
    comprises the places of the undigested matter, the digestive
    fire and the digested matter, of urine and of blood, and the of
    the heart, \sanengdev{uṇḍuka}{intestines} and lungs”
    (\dev{sthānānyāmāgnipakvānāṃ mūtrasya rudhirasya ca/ hṛduṇḍukaḥ
    phuphphusaśca koṣṭha ityabhidhīyate}, \SS\
    \Su{4.2.12--13}{409}).   The occurrence of \dev{uṇḍuka} in a
    list together with \dev{antrāṇi} at \SS\ \Su{3.5.5}{364} would
    suggest that the former was not just “intestines,” but this
    passage does not occur in the Nepalese version of the \SS, so
    ancient thinking on this subject was evidently in flux. See
    further references in \cite[177--178]{das-2003} and
    \volcite{1}[109, 241]{josi-maha}.} It is attached all around the
    liver and the \se{koṣṭha}{abdominal cavity} as well as the
    intestines.\footnote{\Dalhana{3.4.17}{356} discussed this
        passage in detail, citing Gayadāsa, Vṛddhavāgbhaṭa and Caraka.
        Each of these authorities had different interpretations of the
        passage.  For reference, see the diagram of the peritoneum in
        \cite[fig.\,3]{isaz-2018}}
\end{sloka}
 
 
 \item[18] 
 
The sixth, called \se{pittadharā}{bile-support}, is the one that
\se{\root pac}{cooks} the fourfold items known as “eaten, drunk,
chewed and licked".\footnote{\label{fourfoods}Ḍalhaṇa on
    \Su{1.14.3}{59} said that the four types of food are those that can be
    drunk, licked, eaten and chewed(\dev{peyalehyabhojyabhakṣya}). The
    main text of the \CS\ is explicit about these categories at
    \Ca{4.3.4(1)}{308}: \dev{pānāśanabhakṣyalehya/} “things drunk, eaten,
    chewed or licked.” \citet{yagi-1994} discussed the distinction between
    \dev{bhakṣya} and \dev{bhojya}; for further Indological background on
    foods, see the studies by \citet{oliv-1995,oliv-2001} and the classic
    reference works by \citet{acha-1994,acha-1998}. The long, final
    adhyāya of the \SS's \emph{Sūtrasthāna} (ch.\,46) amounts to a
    distinct treatise on food in āyurveda.}
 
 \item[19]
 
 \begin{sloka}
   Whatever item eaten, drunk, chewed and licked, on reaching the
\se{koṣṭha}{innards}, \se{\root jṝ}{decomposes}, being dried in due
time by the heat of bile.
 \end{sloka}
 
\item[20]

The seventh, called \se{śukradharā}{semen-bearer}, is the one that
pervades the whole body of all beings.\footnote{In this case, the
    \dev{kalā} does not correspond in any obvious way to an entity
    identifiable in contemporary anatomy.}

\item[21]

There is a verse on this:
\begin{sloka}   
    The best physician regards the semen in men's bodies just like
ghee in milk or jaggery in cane-juice.\footnote{The metaphors
    point to the essential components in substances that are initially
    invisible but later produce a solid product.}
\end{sloka}

\subsection{Pathways in the body}

\item[22]

The \se{srotas}{tube} for urine is located two \se{aṅgula}{finger's
    breadth} to the right side, below the opening of the bladder. A man's
semen is discharged from it.\footnote{Cf.\ receptacle “F” in the
    figure at \cite[236]{wuja-2008a}.
    
    The vulgate text has an additional
    verse at this point.}.

\item[24]

The paths of the \sepl{srotas}{tube} of pregnant women which carry
menstrual blood become obstructed by the \se{garbha}{fetus}. This is
why pregnant women have menstrual blood, but it is not seen. Then, its
downward motion being obstructed, it accumulates above the upper part
and is then released into the \se{udara}{abdomen}. %
The remainder arrives at an even higher part; it enters into the
breasts. %
That is why pregnant women have full, raised
breasts.\footnote{\citet[81--82]{das-2003} discussed the vulgate
    version of this passage, which has significant differences
    including the coining of the term \dev{aparā} that is not present
    in the Nepalese version.}


\subsection{The formation of the fetus's organs}

\item[25]

The liver and the spleen are produced from blood.\footnote{The vulgate
    text has \dev{garbhasya} “of the fetus” in this sentence.}  The
    \se{pupphusa}{lung} comes from the foam of blood.\footnote{For
        references and interpretation of \dev{pupphusa}, see \cite[271,
        n.\,23, 276, n.\,31]{wuja-2003}.  Prof.\ Zysk noted the Greek parallel
        \textgreek{φῦσα} “bellows” (\cite[1700a]{lidd-1901};
        \volcite{2}[307]{KEWA}).  The Greek word is not commonly used in Greek
        medical literature and the Sanskrit form is reduplicated.}  The
        \se{uṇḍuka}{lower intestine} comes from an \se{kiṭṭa}{secretion} of
        the blood.\footnote{On \dev{uṇḍuka}, see note \ref{uṇḍuka} above.}
        
\item[26]
And there is a verse on this.
\begin{sloka}    
What is considered the best \se{prasāda}{nutriment} of blood and
phlegm is cooked by bile, and also cleansed by the wind.\footnote{The
    \dev{prasāda} “nutriment” is one of the two components of humoral
    substances (or sometimes of the \sepl{dhatu}{body tissue}), the other
    being the \dev{mala} “impurity."  See
    \volcites{1}[536]{josi-maha}{}[563 et passim]{das-2003}. Note that
    blood is here treated as a humoral entity (see \cite{meul-1991} for
    discussion). %\SS\ \Su{1.1.24(2)}{6}).
    Ḍalhaṇa reported that the commentator Gayadāsa added the phrase
    “[being cooked] within the heart."}
\end{sloka}

\item[27--28]
\begin{sloka}
A person's intestines are produced from that, as well as the rectum
and bladder. In that context, the child, being stirred and heated like
gold develops a tongue by which he can recognize tastes.\footnote{The
    ungrammatial \dev{dhāmyamāna} meaning “heating, smelting” rather than
    “blowing”is common in alchemical literature; see
    \cites[passim]{hell-2009}[327, n.\,11]{wujad-2025}.
    
    The vulgate has extra verses here that extend the idea of wind
    as the creator of various further parts of the body.} %Edgerton 455b
    % rukma
\end{sloka}
        
\item[31]
        
The heart originates from the pure part of blood and phlegm. It is
spoken of as the seat of consciousness because of being the location
of the \sepl{dhamanī}{pipe} that carry the life-breath.

When it is covered by tamas, all living beings sleep. 

\subsection{On sleep}

\item[33]

But they teach that the sleep \se{vaiṣṇavī}{related to Viṣṇu} is
\se{pāpman}{injurious}.\footnote{Ḍalhaṇa glossed “related to Viṣṇu" as
    meaning \se{māyā}{divine illusion} (\dev{viṣṇoriyaṃ vaiṣṇavī
    māyeva/}).  The allusion is to Viṣṇu's period of sleep, although it is
    not known why this should be called injurious or wicked.} %
    By nature, this sleep caresses all living beings.
    
In that context, when phlegm, which is dominated by tamas, reaches the
\sepl{srotas}{pipe}\label{sanjnavahasrotas} that convey consciousness,
it is the sleep called “\se{tāmasī}{of tamas}".\footnote{On the “pipes
    that convey consciousness” see references at
    \volcite{1}[868]{josi-maha}.  The conduit is usually a \dev{srotas}
    (\Su{6.61.8}{800}) but exceptionally a \dev{nāḍī} (\SS\
    \Su{6.46.6}{739}), in the context of \se{unmāda}{insanity}.} At the
    time of dissolution one does not wake up from it.  For those who have
    a lot of tamas, it happens during days as well as nights.  For those
    who have a lot of rajas it happens for no reason.\footnote{Ḍalhaṇa
        glossed “for no reason” as “at any time” 
        (\dev{aniyatakāle}); that kind
        of sleep is called “\se{tāmasī}{of tamas}." ).} For those who have a
        lot of sattva, it happens at midnight.  For those who have diminished
        phlegm and abundant wind, it is \se{vaikārikī}{subject to
            modification} because of the damage to the mind.\q{Look more at 
            Ḍalhaṇa.}
            % think a bit more
        % about this.

\item[34]

And there are verses on this:

\begin{sloka}
O Suśruta, the heart is said to be the location of the
consciousness for living beings.  And when it is overpowered by tamas, then 
sleep enters living beings.
\end{sloka}

\item[35]

\begin{sloka}
Tamas is said to be cause of sleep; sattva is the cause of waking.
Alternatively, it is in fact \se{svabhāva}{nature} that is generally
said to be the more weighty cause.\footnote{ “Nature” here probably
    refers to the propensity of an individual.}
\end{sloka}

\item[36]

\begin{sloka}
The lord, as an embodied self, dreams about the experiences in earlier
bodies.  With his mind connected with rajas, he grasps good and bad
things.
\end{sloka}

\item[37]

 \begin{sloka}
Since the organs are imperfect as a result of tamas, even though
the embodied self is not sleeping, he is spoken of as being fast
asleep.\footnote{“Organs,” (\dev{karaṇa}) refers to the sense organs
    and the mind, the “{collection of the senses}” as Ḍalhaṇa put it
    (\Dalhana{3.4.37}{359}:  (\dev{indriyagrāma}). 
    
    This and the previous verse seem to describe two kinds of sleep
that an embodied person experiences: normal dreaming sleep, and a
kind of waking sleep in which cognition is limited by tamas.}
 \end{sloka}

\subsubsection{Against sleeping during the day}

\item[38]

Daytime sleep is prohibited during every season except in
\se{grīṣma}{summer}.\footnote{On the shifting terminology of seasons
    in the \SS\ see \cite{ange-2022}.} Even during the prohibited times,
    what is prohibited is recommended for a period of \se{muhūrta}{three
        quarters of an hour} for children, the elderly, someone
    \se{kṣatakṣīṇa}{with a chest injury},\footnote{\label{angamarda}The
        condition is the subject of \CS\ \Ca{6.11}{478--482} and is
        characterised as damage to or in the chest typically caused by
        over-exertion, and typically presenting as cough and wasting.} people
        who have drunk alcohol,\footnote{The vulgate reads “who are always
            drunk.”} women, and those who are exhausted by the action of
            travelling in a vehicle or on foot, those who have not eaten, and
            those with reduced fat, wind, phlegm and chyle.\footnote{The list of
                vulnerable types, “children, the elderly, \ldots” is a standard
                formula that appears in several other locations in the \SS.} %
                For those who are awake at night, in fact, one may
                sleep during the day for half the time he was awake.

But sleeping in the day is really an abnormality. People sleeping at
that time obtain demerit and also irritation of all their humours;
based on this irritation they develop cough, wheezing, a cold,
headache, heaviness of the head, \se{aṅgamarda}{crushing pain in the
    limbs}, \se{arocaka}{loss of appetite} and weak
digestion.\footnote{The vulgate added “fever” to this list of
    illnesses.
    
    On \dev{aṅgamarda}, “crushing pain in the limbs,” the
\volcite{1}[9]{josi-maha} cites passages from the \SS\ and \CS\
describing pain or tingling in the limbs, as if the limbs were
being rubbed, etc.}  And also, in fact, for those who
are awake during the night, the very same \sepl{doṣa}{problem}
manifest for the same cause.

\item[39]

There are verses about this:
\begin{sloka}
    Therefore, one should not be awake at night and should avoid sleeping 
    during the day.  Having recognized that both cause problems, 
    the wise person should practice measured sleep.   
\end{sloka}

\item[40]
\begin{sloka}    In this way, a man may live intensely for a century free from
disease, with a happy mind,  strong and good-looking, neither too
fat or too thin, and with good fortune.\footnote{The vulgate reads
    \dev{vṛṣaḥ} “as a bull” for \dev{bhṛśam} “intensely,” thus
    introducing the idea of virility \citep[359]{vulgate}.}
\end{sloka}

\subsubsection{Sleep disorders}

\item[56] 
\begin{sloka}
 %   \q{see 41} 
 Fainting consists predominantly of bile and tamas;
dizziness arises from rajas, bile and wind; laziness arises from
tamas, wind and phlegm, sleep originates in phlegm and tamas.
\end{sloka}

\item[42]

\begin{sloka}
    Insomnia occurs due to wind, bile, heat of the mind (i.e.,
\se{manastāpa}{worries}), and also due to \se{kṣaya}{wasting}, as well
as injury. It is alleviated by their 
countermeasures.\footnote{\Dalhana{5.4.42}{359} glossed 
“countermeasures” as treatments like massage.}
\end{sloka}

%\item[43]
%
%In insomnia oil massage of the head, body rubs, powder massage of the body 
%and gentle stroking massages are beneficial.
%
%\item[44]
%
%With meals consisting of preparations made from ground rice and wheat flour, 
%prepared  with sugarcane products; meals that are sweet, unctuous and 
%enriched with milk, meat  broths, and similar nourishing items.
%
%\item[45]
%
%At night, it would be appropriate to apply a diet with broths of burrowing 
%animals and birds that scatter grain, as well as products with grapes, of white 
%sugar and sugarcane preparations.
%
%\item[46]
%
%One should arrange soft  and  pleasing beds and seats. In case of 
%sleeplessness ; however , the wise person should also employ  other 
%appropriate 
%measures.

\item[47]

\begin{sloka}
    In the case of excessive sleep, the following are beneficial:
emesis, cleansing therapies, fasting, blood-letting and
\se{manovyākulana}{mental diversions}.
\end{sloka}

\item[48]

\begin{sloka}
    For those afflicted by phlegm, fat, or poison, it is good to stay awake at
night. And sleeping during the day is good 
for those suffering from thirst, sharp pain, hiccups,
indigestion, or diarrhoea.
\end{sloka}

\item[49]

\begin{sloka}
    A person with a sleep disorder is here defined as having
\se{tandrā}{lassitude} if they are unaware of the objects of the sense
organs, if they yawn, feel heavy and exhausted.\footnote{The vulgate adds 
    six verses here.}.
\end{sloka}

%\item[50]
%
%Taking in a single deep breath, twisting the body, and opening the
%mouth wide, one then releases the air while tears well in the eyes;
%this action is known as a yawn.
%
%\item[51]
%
%That bodily fatigue which arises in the absence of physical exertion, is
%intense, and occurs without shortness of breath, should be understood
%as \se{klama}{languor}; it is a condition that obstructs the functioning
%of the senses.
%
%\item[52]
%
%Indulgence in the pleasureable contact, the fickleness of hating pain,
%and a lack of exertion toward activities even when one is capable are
%defined as \se{ālasya}{indolence}.
%
%\item[53]
%
%When food causes nausea but does not emerge, though it is agitated by
%excessive salivation and spitting, and the individual feels a
%distressing pressure in the heart—that condition should be identified
%as \se{utkleśa}{nausea}.
%
%\item[54]
%
%Sweetness in the mouth, lethargy, a feeling of constriction in the
%heart, dizziness, and having  no desire for food—such
%should be identified as the signs that he has \se{glāni}{exhaustion}.
%
%\item[55]
%
%The man who regards his body as merely wrapped in a damp skin, and
%whose head is heavy is defined as having \se{gaurava}{heaviness}.


\subsection{Again on the formation of the fetus's organs}

\item[57]

In fact, the development of the fetus depends on the quality of the
nutritional juice and on the food consumed by the
mother.\footnote{This abrupt change of topic signals a textual disturbance in
    the text's transmission at this point.  Perhaps the preceding section on sleep 
    is an old intrusion.}

\item[60]

An there is a verse on this:
\begin{sloka}
     It is  the opinion of Dhanvantari that
     the \se{dṛṣṭi}{pupil} and the pores never grow.
     These are fixed for human beings.\footnote{Dhanvantari is here apparently 
     an outside authority.}
\end{sloka}

\item[61]

It is an established fact that nails and hair establish their 
\se{prakṛti}{character} 
naturally and then always grow, even when the body decays.


\subsection{On temperaments}

\item[62]

There are three basic \sepl{prakṛti}{temperament}, caused by wind,
bile, and phlegm.\footnote{The vulgate presents seven temperaments
    here, not three.}

\item[63] 

The humour that is in the ascendent when semen and menstrual blood join 
gives rise to the temperament.  Listen to my characterisation of them.


\subsubsection{Wind}

\item[64]

\q{pneumatic?}Among these, the person with a wind temperament is wakeful, adverse to
cold, an unhappy thief, ignoble, greedy, given to music, has split
feet, has rough, sparse hair and moustache, is angry with
difficulties, grinds his teeth, is weak and has a short
life.\footnote{\Dalhana{3.4.64}{361} noted that the grinding of the
    teeth happens while the person is asleep.}
 
\item [65--66]

On this:
\begin{sloka}
    A person who has a wind temperament is unsteady,  unreliable in friendship, ungrateful, thin and chapped, covered with \sepl{dhamanī}{vein}, 
    garrulous, he wanders about with rapid gait, and his \se{ātman}{self} is unsteady.  When he is asleep, he roams about in the clouds.
\end{sloka}
\begin{sloka}
    His opinions are unsettled, his thoughts are flighty.  He is
slow to get teeth, wealth or friends.\footnote{This slightly odd
    expression reads like a worldly-wise saw, though it is not in \cite{jaco-handful}; the vulgate reads
    differently.} And he babbles unconnectedly about whatever.
\end{sloka}
 
 \subsubsection{Bile}
 
 \item[68]
 
\q{Choleric?}Someone of a bilious temperament is sweaty, can tolerate the cold, ill-smelling, \se{pīta}{jaundiced} or \se{kāla}{black}.\footnote{This could be a notable reference to yellow and black bile.} His limbs are flaccid, his nails, eyes, palate, tongue, lips, soles and palms have a coppery colour.  He is unlucky, covered in wrinkles and grey hair, he eats a lot, dislikes heat, is habitually quick to anger and quick to calm down, and he has a medium span of life. 
    
\item[69--70]

And on this:
\begin{sloka}
    He is an intelligent, sharp-minded speaker who defeats
opponents.\footnote{The term \dev{nigṛhyavaktā} almost certainly
    refers to the rules of formal debate as described in \CS's
    Vimānasthāna~8  (see, e.g., \volcite{2}[71, 129]{pret-1991} on
    debate.  See also \emph{ibid}.\ 2, 129--132 on
    \dev{nigrahasthānam} and \emph{ibid}.\ 3, 136--138, on the related
    term \dev{vigṛhyasaṃbhāṣā}). Cf.\ the use of the term in the
    \emph{Mahāyānasūtralaṅkāra}, 20 \& 21: 53,
    \parencites[186]{levi-1907}[181]{dvar-1985}; translated, “Tu
    parles en censeur dans les assembl\'ees” by \citet[303]{levi-1907}}
    He is brilliant in conflicts, with unyeilding valiance. Furthermore, while he
    is sleeping he sees \gls{kanaka}, \gls{palāśa}, and
    \gls{karṇikāra}, as well as fire, lightning and meteors.
\end{sloka}
\begin{sloka}
    He does not bow down out of fear, but he is gentle toward those who bow to him.  And to those who bow he is gentle and generous. 
    In this world, his facial expressions and gait always appear distressed. Such a person is one whose temperament is made of bile.\footnote{The vulgate here has a verse, \Su{3.4.71}{361}, identifying bilious types according to the features of various animals, etc. The editors of the vulgate noted that this verse does not appear in manuscripts (\emph{idem}, n.\,2).}
\end{sloka}


 \subsubsection{Phlegm}
 
 \item[72]
 
 A person of phlegmatic temperament is characterized by a complexion that is any one of \gls{dūrvā}, \gls{indīvara}, \gls{nistriṃśapatra}, 
 fresh \gls{ariṣṭaka}, or \gls{śara} stalk.\footnote{Strangely, mostly greenish hues.}  Such a person is  lucky, good-looking, and fond of sweet things.  He is appreciative, steadfast, tolerant, free from strong desire, strong, 
     retaining for a long time, firm in enmity, and long-lived.
 
 
 \item[73]
 
\begin{sloka}
    He has an shiny body, firm, well-proportioned, and beautiful limbs,
he is handsome, and has a voice like thunder, kettle drums, or a
lion. Furthermore, when asleep, he dreams of lovely lakes full of
lotuses, swans, and \glspl{cakravāka}.\footnote{This verse and v.\,69 
    share a sentence structure that is clarified by the testimonia in \AS\ and \Ah\ cited in the critical edition's apparatus.}
\end{sloka}

\item[75]
 
\begin{sloka}
     A person of phlegmatic temperament has a firm  grasp of the shastras,
he is a steady friend, and after careful calculation he donates much.
The words in his sentences are always perfect and he respects his
elders.
\end{sloka}

\item[77.add]

\begin{sloka}
    When, in a person’s constitution, the form of a pair of humours is
seen, one should recognize that as being a “\se{saṃsarga}{combined
    constitution}.” Also, a “\se{saṃsarga}{combined constitution}”
which is threefold.\footnote{The point here is that a “combined”
    constitution formed by two humours can be of three types:
    wind-bile, wind-phlegm, and bile-phlegm.  Another interpretation
    could be that there is another kind of combination that is
    threefold. The more usual term for a pathological combination of
    three humours is \dev{sannipāta}.}
\end{sloka}

\item[78]

\begin{sloka}
    Neither the irritation nor alteration, nor depletion of constitutions
occur naturally.  But they are observed in someone whose life is over.
\end{sloka}

\item[79]

\begin{sloka}
A poisonous insect does not succumb to its poison.  In the same way, the body is not afflicted by its constitutions, because it is born from them.
\end{sloka}
        
%\item[50]
%\item[51]
%\item[52]
%\item[53]
%\item[54]
%\item[55]
%
%\item[56]  
%\item[58]
%\item[59]


\item[80]

Some say that in this system % iha
the constitution of men is derived from the elements. Those three that have been mentioned arise due to wind, fire, and water. If a person is earthy, he has a steady and large body and is patient. %
And the ethereal person % nābhasaḥ ?? 
is pure, long-lived and has large apertures. % khair mmahadbhiḥ
\newpage

\subsection{Typology of Bodies}

\subsubsection{Seven sattvic bodies}

\item[81]

The characteristics of someone having a Brahmā-type body are 
purity, belief in the next world, % āstikya
reciting the Vedas, %  abhyāso vedeṣu
honouring one's elders, % gurupūjanam
hospitality, % priyātithitvam
and ritual worship. % ijyā


% got to here 3 July 2026


\item[82]

The characteristic of the Great Indra body (\emph{mahendra}) is heroism, authority, great fortune, always knowing \se{śāstra}{authoritative teachings}, and supporting servants.  

\item[83]

The characteristic of the Varuṇa body is endurance, being used to the
cold, a tawny complexion, yellow hair, and loving
water.\footnote{Varuṇa is the deity presiding over waters, amongst
    other things.}

% got to here 2026-07-17

\item[84]

The characteristics of the Kubera-type body are impartiality,
endurance, the acquisition and accumulation of wealth, and great
procreative power.

\item[85]

The characteristics of the Gandharva-type body are a fondness for
perfumes and garlands,  a love for dance and music, and an inclination
for roaming about.

\item[86]

The person of a Yama nature behaves correctly, is firm in
undertakings, is resolute, has a good memory, is pure, and free from
passion, hatred, fear, and ignorance.

\item[87]

One knows a person of a sagely nature to be somebody who is devoted to celibacy,
vows,\footnote{Or “the vow of celibacy.”} mantra recitation, fire rituals, and
study; they have both wisdom and particular knowledge.\footnote{Elsewhere,
\Dalhana{1.35.5--6}{149} distinguished \dev{jñāna} and \dev{vijñāna} as
\dev{tattvāvabodhaḥ} and \dev{citrādikarmakauśalyam} “insight into reality” and
“expertise in activities like painting.”}

\item[88--89.add]

These are the seven sattvic bodies.  
\begin{table}
    \centering
    \caption{Sequence of rajasic body types}
    \label{rajasic-bodies}
    \bigskip
%    \begin{tabular}{>{\it}lcc}
%        & \emph{Nepalese} & \emph{vulgate} \\
%        \toprule
%        āsura & 1 & 1 \\   
%        rākṣasa & 2 & 4 \\
%        paiśāca & 3 & 5 \\
%        sarpasattva & 4 & 2 \\
%        pretasattva & 5 & 6 \\
%        śākuna  & 6 & 3 \\
%        \bottomrule
%    \end{tabular}
\small

\begin{tikzpicture}[
    >=Latex,
    %thick,
    lbl/.style={font=\itshape},
    hdr/.style={font=\rmfamily}
    ]
% --- column headers ---
\node[hdr] (H1) at (1,0.8)   {Nepalese};
\node[hdr] (H2) at (4.5,0.8) {vulgate};
%\draw (-1.3,0.32) -- (4.3,0.32);


% --- left column: Nepalese sequence, top to bottom ---
\node[lbl,anchor=west] (L1) at (0,0.0)   {1. āsura};
\node[lbl,anchor=west] (L2) at (0,-0.8) {2. rākṣasa};
\node[lbl,anchor=west] (L3) at (0,-1.6) {3. paiśāca};
\node[lbl,anchor=west] (L4) at (0,-2.4) {4. sarpasattva};
\node[lbl,anchor=west] (L5) at (0,-3.2) {5. pretasattva};
\node[lbl,anchor=west] (L6) at (0,-4.0) {6. śākuna};

% --- right column: vulgate sequence, top to bottom (note: no "6" in source) ---
\node[lbl,anchor=west] (R1) at (4,0.0)   {1. āsura};
\node[lbl,anchor=west] (R2) at (4,-0.8) {2. sarpasattva};
\node[lbl,anchor=west] (R3) at (4,-1.6) {3. śākuna};
\node[lbl,anchor=west] (R4) at (4,-2.4) {4. rākṣasa};
\node[lbl,anchor=west] (R5) at (4,-3.2) {5. paiśāca};
\node[lbl,anchor=west] (R6) at (4,-4.0) {6. pretasattva};

% --- curved arrows, left to right, one per entity ---
\draw[->] (L1.east) to[out=0,in=180] (R1.west); % asura
\draw[->] (L2.east) to[out=0,in=180] (R4.west); % raksasa
\draw[->] (L3.east) to[out=0,in=180] (R5.west); % paisaca
\draw[->] (L4.east) to[out=0,in=180] (R2.west); % sarpasattva
\draw[->] (L5.east) to[out=0,in=180] (R6.west); % pretasattva
\draw[->] (L6.east) to[out=0,in=180] (R3.west); % sakuna

\end{tikzpicture}
\end{table}

\subsubsection{Six rajasic bodies}


Learn from me about the rajasic ones.\footnote{As above, the text
    describes bodily charactersistics as personality traits and
    behaviours. The lists of qualifiers are in the accusative
    case, as objects of \dev{nibodha} “learn from me.”  The sequence of body-types in the
    Nepalese version accords with the commentary of Gayadāsa but differs
    from the vulgate; see Table~\ref{rajasic-bodies}.}

\item[88cd]
%
The Asura  \se{sattva}{nature} is like this: 
magisterial, wild, heroic, wrathful, and jealous. 
He eats alone, he is deceitful.
%


\item[89.add, 92ab]

The characteristic of the \se{rākṣasa}{goblin} body %
is that he is fixated on one view, is wild, and he speaks enviously
and unrighteously.  And he also brags excessively.

\item[92cd--93ab]

The qualities of the \se{paiśāca}{imp} include recklessness, coarseness,
and a tendency to impetuousness.  He is a lecher, he is without shame. He is 
sharp, full of travail, angry, and essentially timorous. 


\newpage
\begin{tt}
    \raggedright
    \bigskip
    \marginpar{Draft tr.\ from here}
    
%\item[90]  cf. infra
%
%\item[91] cf. infra




\item[90] (cf.supra?)

(A man) given to pleasure and food, fickle, and of snake-like nature — they consider such a man  to be ungenerous, lazy, of bad character, and untrustworthy.

\item[94]

They consider a greedy and  ungenerous man to have the nature of a ghost.

\item[91]  (cf. supra?)

One whose desires are awakened, or who avoids enjoyment,
who constantly eats, and who indeed (is) intolerant, has no fixed abode— this is called a bird-like nature. These six are Rājasic traits; but now learn from me about the Tamas-type.

\subsection{Three tamasic bodies}

\item[95]

Poor understanding, slowness, excessive sleep, habitual indulgence in sex, incapacity (to take an initiative), and sorrows — these are to be recognized as beastly qualities.

\item[96]

Instability, foolishness, cowardice, fondness for water, and mutual jostling — these are qualities entirely fish-like in nature.

\item[97]

Lacking spirit and strength of limb, devoted solely to food, without cherished aims or desires  — that would be a man with the nature of a tree).

\item[98]

Thus these threefold constitutions, beginning with the rājasic and subsequently as discussed, have been described. Having recognized one’s bodily constitution, one should act accordingly, so it is said. 

\item[99]

Thus ends the Śārīra, the fourth chapter.

\end{tt}
\end{translation}